% frontmatter/abstract.tex — written by hand for this paper.
% Included right after \maketitle by the generated main.tex.

\begin{abstract}
The term \emph{data} is among the most frequently used and least precisely
defined concepts in computing. Standard references converge on an informal
characterization of data as ``raw facts,'' but none specify the formal
structure that makes a collection of facts addressable, comparable, or
processable by a machine. This representational ambiguity propagates into
system design as schema drift, serialization defects, and a class of
type-confusion vulnerabilities that arise when a fixed byte sequence is
silently reinterpreted under an unstated indexing or typing assumption.
This paper makes the case for treating data formally as a
\emph{function from an index set to a value set}, sharpens this definition
with explicit well-definedness conditions, and positions it against three
existing frameworks: Floridi's General Definition of Information, Zins'
faceted classification, and the DIKW hierarchy. We derive four corollaries
connecting this definition directly to recurring failure modes in software
engineering---implicit schema coupling, byte-level reinterpretation,
unsound equality checks, and ambiguous serialization contracts---with
worked examples for each.
\end{abstract}

\paragraph{Keywords} data representation; formal definition; information theory;
type safety; serialization; systems design; data semantics.